% Source: MQ14b_v1.html
\documentclass[11pt]{article}
\usepackage[margin=1in]{geometry}
\usepackage{amsmath}
\usepackage{amssymb}
\usepackage{siunitx}
\usepackage{enumitem}

\title{MQ14b: Oscillations}
\author{}
\date{}

\begin{document}
\maketitle

\section*{MQ14b: Oscillations}

\begin{enumerate}[leftmargin=*,label=\textbf{Question \arabic*.}]
\item \hfill \textit{1 pts}

A simple harmonic oscillator consists of a block of an unknown mass in kilograms and spring with spring constant 689 N/m on a frictionless horizontal surface. The block is moved a distance 1.58 mfrom equilibrium and released from rest. If the block's maximum speed at the equilibrium point is 14.8 m/s, calculate the mass of the block in kg.

\item \hfill \textit{1 pts}

Suppose a mass-spring system has spring constant 665 N/m and mass 5 kg. The mass is displaced in the positive x-direction by 0.2 m and set in motion at 2.1 m/s in the negative x-direction. Find the amplitude \emph{A} of the motion in m.

\item \hfill \textit{1 pts}

A mass-spring system in simple harmonic motion has an energy of 7.98 J and an amplitude of 14 cm. Find the maximum force that the spring exerts on the mass. (Be careful with units!)



\emph{F }= \rule{2cm}{0.4pt} N

\item \hfill \textit{1 pts}

Calculate the length of a simple pendulum on Earth if its angular frequency is 2.74 rad/s.



 \emph{L} = \rule{2cm}{0.4pt} m

\item \hfill \textit{1 pts}

Suppose an irregular chunk of matter of mass 13.6 kg has a moment of inertia 11.18 kgᐧm$^{2}$ about its pivot located 1.74 m away from its center of mass. Calculate the angular frequency in rad/s of small oscillations.

\end{enumerate}

\end{document}
